% Part: turing-machines
% Chapter: undecidability
% Section: trakhtenbrot

\documentclass[../../../include/open-logic-section]{subfiles}

\begin{document}

\olfileid[pt]{tur}{und}{tra}
\olsection{Teorema de Trakhtenbrot}

\begin{explain}
  Na \olref[rep]{sec} definimos as !!{sentence}s $!T(M,w)$ e~$!E(M,w)$
  para uma máquina de Turing~$M$ e uma cadeia de entrada~$w$. Em
  seguida, mostramos em \olref[ver]{lem:valid-if-halt} e
  \olref[ver]{lem:halt-if-valid} que $!T(M,w) \lif !E(M,w)$ é válida
  se, e somente se, $M$, iniciada na entrada~$w$, eventualmente para.
  Como o Problema da Parada é indecidível, isso implica que a validade
  e a satisfatibilidade de !!{sentence}s da lógica de primeira ordem
  são indecidíveis
  (\cref{tur:und:uns:thm:decision-prob,tur:und:uns:cor:undecidable-sat}).

  Mas a validade e a satisfatibilidade de sentenças são definidas para
  !!{structure}s arbitrárias, finitas ou infinitas. Você poderia
  suspeitar que é mais fácil decidir se !!a{sentence} é satisfatível em
  uma !!{structure} finita (ou válida em todas as !!{structure}s
  finitas). Podemos adaptar a prova da insolubilidade do problema da
  decisão de modo que ela mostre que não é esse o caso.

  Primeiro, se você voltar à prova de
  \olref[ver]{lem:halt-if-valid}, verá que o que fizemos ali foi
  produzir um modelo~$\Struct M$ de~$!T(M,w)$ que descreve exatamente
  o que a máquina~$M$ faz quando iniciada na entrada~$w$.  O domínio
  daquele modelo era~$\Nat$, isto é, infinito. Mas se $M$ de fato para
  na entrada~$w$, podemos construir um modelo finito~$\Struct M'$ da
  mesma maneira. Suponha que $M$ iniciada na entrada~$w$ pare após~$k$
  passos. Tome como domínio~$\Domain{M'}$ o conjunto $\{0, \dots, n\}$,
  em que $n$ é o maior entre~$k$ e o comprimento de~$w$, e seja
  \[
    \Assign{\prime}{M'}(x) = 
    \begin{cases}
      x + 1 &\text{se $x < n$}\\
      n &\text{caso contrário,}
    \end{cases}
  \]
  e $\tuple{x,y} \in \Assign{<}{M'}$ se, e somente se, $x < y$ ou $x =
  y = n$. Quanto ao resto, $\Struct{M'}$ é definida tal como~$\Struct{M}$.
  Pela definição de~$\Struct{M'}$, tal como na prova de
  \olref[ver]{lem:halt-if-valid}, $\Sat{M'}{!T(M,w)}$.  E como
  supusemos que $M$ para na entrada~$w$, $\Sat{M'}{!E(M,w)}$. Assim,
  $\Struct{M'}$ é um modelo finito de~$!T(M,w) \land !E(M,w)$ (note
  que substituímos $\lif$ por~$\land$).
  
  Estamos a meio caminho de uma prova: mostramos que se $M$ para na
  entrada~$w$, então $!T(M,w) \land !E(M,w)$ tem um modelo finito.
  Infelizmente, a recíproca disso não vale, isto é, há máquinas de
  Turing que não param em alguma entrada~$w$, mas $!T(M,w) \land
  !E(M,w)$ ainda assim tem um modelo finito. Por exemplo, considere a
  máquina~$M$ com o único estado $q_0$ e a instrução
  $\delta(q_0,\TMblank) = \tuple{q_0,\TMblank,\TMstay}$. Iniciada na
  entrada vazia~$w = \emptyseq$, essa máquina nunca para: está em um
  laço infinito, mas não altera a fita nem move o cabeçote. Todas as
  configurações são as mesmas (mesmo estado, mesma posição do
  cabeçote, mesmo conteúdo da fita). Podemos definir uma !!{structure}
  finita~$\Struct{M''}$ que satisfaz $!T(M,\emptyseq) \land
  !E(M,\emptyseq)$ (exercício).  Podemos, no entanto,
  modificar~$!T(M,w)$ de modo adequado para que tais !!{structure}s
  sejam descartadas.

  \begin{prob}
    Seja $M$ uma máquina de Turing com o único estado $q_0$ e a única
    instrução $\delta(q_0,\TMblank) = \tuple{q,\TMblank,\TMstay}$.
    Seja $\Domain{M''} = \{0, 1, 2\}$, $\Assign{\prime}{M''}(0) =
    \Assign{\prime}{M'}(1) = 1$ e $\Assign{\prime}{M''}(2) = 2$, e
    $\Assign{<}{M''} = \{\tuple{0,1}, \tuple{1,1}, \tuple{2,2}\}$.
    Defina $\Assign{\Obj Q_{q_0}}{M''}$, $\Assign{\Obj
    S_{\TMblank}}{M''}$ e $\Assign{\Obj S_{\TMendtape}}{M''}$ de modo
    que $!T(M,\emptyseq)$ e $!E(M,\emptyseq)$ se tornem verdadeiras e
    explique por que o são. Dica: Observe que $\delta(q_0, \TMendtape)$
    é indefinida. Garanta que
    \begin{align*}
    & \Obj Q_{q_0}(\num{1}, \num{n}) \land \Obj S_{\TMendtape}(\num{0}, \num{n})
      \land 
      \lforall[x][(\num{0} < x \lif \Obj S_\TMblank(x, \num{n}))] \text{\quad para todo $n \in \Nat$}\\
    & \lexists[y][(\Obj Q_{q_0}(\num{0}, y) \land \Obj S_{\TMendtape}(\num{0}, y))]
    \end{align*}
    são ambas verdadeiras em~$\Struct{M''}$.
  \end{prob}
\end{explain}

Considere as !!{sentence}s que descrevem a operação da máquina de
Turing~$M$ na entrada $w = \sigma_{i_1}\dots\sigma_{i_k}$:
\begin{enumerate}
\item Axiomas que descrevem números e~$<$ (tal como na definição
de~$!T(M,w)$ na \olref[rep]{sec}).
\item Axiomas que descrevem a configuração de entrada: tal como na
definição de~$!T(M,w)$.
\item Axiomas que descrevem a transição de uma configuração para a
  seguinte:

Para o que segue, seja $!A(x, y)$ como antes, e seja
\[
  !B(y) \ident \lforall[x][(x < y \lif \eq/[x][y])].
\]
\begin{enumerate}
\item \ollabel{rep-right} Para toda instrução $\delta(q_i, \sigma) =
  \tuple{q_j, \sigma', \TMright}$, a !!{sentence}:
\begin{align*}
& \lforall[x][\lforall[y][(
   (\Obj Q_{q_i}(x, y) \land \Obj S_{\sigma}(x, y)) \lif {}]] \\
&\qquad   (\Obj Q_{q_j}(x', y') \land \Obj S_{\sigma'}(x, y') \land
!A(x, y) \land !B(y')))
\end{align*}
\item \ollabel{rep-left} Para toda instrução $\delta(q_i, \sigma) =
  \tuple{q_j, \sigma', \TMleft}$, a !!{sentence}
\begin{align*}
& \lforall[x][\lforall[y][
    ((\Obj Q_{q_i}(x', y) \land \Obj S_{\sigma}(x', y)) \lif {}]]\\
& \qquad   (\Obj Q_{q_j}(x, y') \land \Obj S_{\sigma'}(x', y') \land
!A(x, y))) \land {}\\
& \lforall[y][((\Obj Q_{q_i}(\num{0}, y) \land \Obj S_{\sigma}(\num{0},
    y)) \lif {}]\\
& \qquad (\Obj Q_{q_j}(\num{0}, y') \land \Obj S_{\sigma'}(\num{0},
  y') \land !A(\num{0}, y) \land !B(y')))
\end{align*}
\item \ollabel{rep-stay} Para toda instrução $\delta(q_i, \sigma) =
  \tuple{q_j, \sigma', \TMstay}$, a !!{sentence}:
\begin{align*}
& \lforall[x][\lforall[y][(
   (\Obj Q_{q_i}(x, y) \land \Obj S_{\sigma}(x, y)) \lif {}]] \\
&\qquad (\Obj Q_{q_j}(x, y') \land \Obj S_{\sigma'}(x, y') \land
!A(x, y) \land !B(y')))
\end{align*}
\end{enumerate}
Como você pode ver, as !!{sentence}s que descrevem as transições de~$M$ 
são as mesmas que a !!{sentence} correspondente
em~$!T(M,w)$, exceto que adicionamos $!B(y')$ ao final. $!B(y')$~garante
que o número $y'$ da ``próxima'' configuração é diferente de
todos os números anteriores $\Obj 0$, $\Obj 0'$, \dots.
% \item Sentences that express consistency conditions:
% \begin{enumerate}
%   \item\ollabel{rep-con-symb}%
%   !!^a{sentence} that says that no square can have more than one
%   symbol on it at any given time:
%   \[\lforall[x][\lforall[y][(\Obj S_{\sigma}(x, y) \lif \lnot \Obj
%   S_{\sigma'}(x, y))]],\]
%   for every pair $\sigma \neq \sigma'$ of tape symbols of~$T$.
%   \item\ollabel{rep-con-state}%
%   !!^a{sentence} that says that $M$ cannot be in more than one state
%   at any given time:
%   \[\lforall[x][\lforall[y][(\Obj Q_{q}(x, y) \lif
%   \lforall[z][\lnot \Obj
%   Q_{q'}(z, y))]],\]
%   for every pair $q \neq q'$ of states of~$T$.
%   \item\ollabel{rep-con-tape}%
%   !!^a{sentence} that says that $M$ cannot be on more than one tape
%   square at any given time:
%   \[\lforall[x][\lforall[y][(\Obj Q_{q}(x, y) \lif
%   \lforall[z][\eq/[z][y]\lif \lnot \Obj
%   Q_{q'}(z, y))]],\]
%   for every pair $q$, $q'$ of states of~$T$.
% \end{enumerate}
\end{enumerate}
Seja $!T'(M, w)$ a conjunção de todas as !!{sentence}s acima para a
máquina de Turing~$M$ e a entrada~$w$.

\begin{lem}\ollabel{lem:halts-sat}
  Se $M$ iniciada na entrada~$w$ para, então $!T'(M,w) \land !E(M,w)$
  tem um modelo finito.
\end{lem}

\begin{proof}
  Seja $\Struct{M'}$ como na prova de
  \olref[ver]{lem:halt-if-valid}, exceto que
  \begin{align*}
    \Domain{M'} & = \{0, \dots, n\},\\
    \Assign{\prime}{M'}(x) & = 
    \begin{cases}
      x + 1 &\text{se $x < n$}\\
      n &\text{caso contrário,}
    \end{cases} \\
    \tuple{x,y} \in \Assign{<}{M'} &\text{sse $x < y$ ou $x = y = n$,}
  \end{align*}
  em que $n = \max(k,\len{w})$ e $k$~é o menor número tal que
  $M$ iniciada na entrada~$w$ parou após~$k$ passos. Deixamos como
  exercício a verificação de que $\Sat{M'}{!T'(M,w) \land E(M,w)}$.
\end{proof}

\begin{prob}
  Complete a prova de \olref[tur][und][tra]{lem:halts-sat}
  provando que $\Sat{M'}{!T(M,w) \land E(M,w)}$.
\end{prob}

\begin{lem}\ollabel{lem:sat-halts}
  Se $!T'(M,w) \land !E(M,w)$ tem um modelo finito, então $M$ iniciada
  na entrada~$w$ para.
\end{lem}

\begin{proof}
  Mostramos a contrapositiva. Suponha que $M$ iniciada em~$w$ não
  pare. Se $!T'(M,w) \land !E(M,w)$ não tem modelo algum, terminamos.
  Então suponha que $\Struct{M}$ seja um modelo de~$!T(M,w) \land !E(M, w)$.
  Temos de mostrar que ele não pode ser finito.
  
  Podemos provar, tal como em \olref[ver]{lem:config}, que se~$M$,
  iniciada na entrada~$w$, não parou após~$n$ passos, então $!T'(M,w)
  \Entails !C(M, w, n) \land !B(\num{n})$. Como $M$ iniciada na
  entrada~$w$ não para, $!T'(M,w) \Entails !C(M, w, n) \land
  !B(\num{n})$ para todo~$n \in \Nat$. Note que, por
  \olref[rep]{prop:mlessk}, $!T'(M,w) \Entails \num{k} < \num{n}$ para
  todo~$k < n$. Além disso, $!B(\num{n}) \Entails \num{k} < \num{n} \lif
  \eq/[\num{k}][\num{n}]$. Assim, $\Sat{M}{\eq/[\num{k}][\num{n}]}$ para
  todo~$k < n$, isto é, os infinitos termos~$\num{k}$ devem ter todos
  valores diferentes em~$\Struct{M}$. Mas isso requer que
  $\Domain{M}$ seja infinito, de modo que $\Struct{M}$ não pode ser um
  modelo finito de~$!T'(M,w) \land !E(M, w)$.
\end{proof}

\begin{prob}
  Complete a prova de \olref[tur][und][tra]{lem:sat-halts} provando
  que se~$M$, iniciada na entrada~$w$, não parou após~$n$ passos,
  então $!T'(M,w) \Entails !B(\num{n})$.
\end{prob}

\begin{thm}[Teorema de Trakhtenbrot]
  \ollabel{thm:trakhtenbrodt}
  É indecidível se uma !!{sentence} arbitrária da lógica de primeira
  ordem tem um modelo finito (isto é, é finitamente satisfatível).
\end{thm}

\begin{proof}
  Suponha que houvesse uma máquina de Turing~$F$ que decide o problema
  da satisfatibilidade finita. Então, dada qualquer máquina de
  Turing~$M$ e entrada~$w$, poderíamos computar a sentença~$!T'(M,w)
  \land !E(M,w)$ e usar~$F$ para decidir se ela tem um modelo finito.
  Por \cref{tur:und:tra:lem:halts-sat,tur:und:tra:lem:sat-halts}, ela
  tem se, e somente se, $M$ iniciada na entrada~$w$ para. Assim,
  poderíamos usar $F$ para resolver o problema da parada, que sabemos
  ser insolúvel.
\end{proof}

\begin{cor}\ollabel{cor:fproof-incomp}%
  Não pode haver um sistema de !!{derivation} que seja correto e
  completo para a validade \emph{finita}, isto é, um sistema de
  !!{derivation} que tenha $\Proves !B$ se, e somente se,
  $\Sat{M}{!B}$ para toda !!{structure} finita~$\Struct{M}$.
\end{cor}

\begin{proof}
  Exercício.
\end{proof}

\begin{prob}
  Prove \olref[tur][und][tra]{cor:fproof-incomp}. Observe que $!B$ é
  satisfeita em toda !!{structure} finita se, e somente se, $\lnot !B$
  não é finitamente satisfatível.  Explique por que a satisfatibilidade
  finita é semidecidível no sentido de
  \olref[tur][und][uns]{thm:valid-ce}. Use isso para argumentar que se
  houvesse um sistema de !!{derivation} para a validade finita, então
  a satisfatibilidade finita seria decidível.
\end{prob}


\end{document}
